\subsection{多项式根中的对称函数}

考虑一个次数为 \(n\) 、系数在 \(A\) 中的首一多项式：

\[
f = \mathbf {X} ^ {n} + a _ {1} \mathbf {X} ^ {n - 1} + \dots + a _ {n - 1} \mathbf {X} + a _ {n}.
\]

我们定义一个结合的、交换的且保单位的 A-代数 \(\mathbf{E}_f\) ，它以 \(x_{1},\ldots ,x_{n}\) 为生成元，关系式为

\[
\sum_ {i _ {1} <   \dots <   i _ {k}} x _ {i _ {1}} \dots x _ {i _ {k}} = (- 1) ^ {k} a _ {k} \quad (1 \leqslant k \leqslant n).\tag{24}
\]

更精确地说，我们有

\[
\mathrm{E} _ {f} = \mathrm{A} [ \mathrm{X} _ {1}, \dots , \mathrm{X} _ {n} ] / \mathfrak {a}
\]

其中理想 \(\mathfrak{a}\) 由多项式 \(s_k + (-1)^{k+1} a_k\) （ \(1 \leqslant k \leqslant n\) ）生成，而 \(x_i\) 是 \(\mathbf{X}_i\) 模 \(\mathfrak{a}\) 的剩余类（ \(1 \leqslant i \leqslant n\) ）。关系式 (24) 也等价于 \(f(\mathbf{X}) = \prod_{i=1}^{n} (\mathbf{X} - x_i)\) 。当有任何歧义风险时，我们写作 \(x_{1,f}, \ldots, x_{n,f}\) 以代替 \(x_1, \ldots, x_n\) 。

命题4. --- 设 B 为一个交换环， \(\rho\) 为从 A 到 B 的一个同态，以及 \(\xi_1, \ldots, \xi_n\) 为 B 的元素。假设关系 \( {}^{\rho}f(X) = \prod_{i=1}^{n}(X - \xi_i)\) 在 B{[}X{]} 中成立。则存在一个且仅一个环同态 \(u: E_f \to B\) ，使得 \(\rho(a) = u(a, 1)\) 对所有 \(a \in A\) 成立，且 \(u(x_i) = \xi_i\) 对 \(1 \leqslant i \leqslant n\) 成立。

我们通过 \(\rho\) 把 B 视为结合的、交换的且保单位的 A-代数。于是关系 \({}^{\rho}f(X) = \prod_{i=1}^{n}(X - \xi_i)\) 也可以写成如下形式

\[
\sum_ {i _ {1} <   \dots <   i _ {k}} \xi_ {i _ {1}} \dots \xi_ {i _ {k}} = (- 1) ^ {k} a _ {k} \cdot 1 \quad (1 \leqslant k \leqslant n)
\]

在 B 中。由于关系式 (24) 给出了 \(\mathbf{E}_f\) 的一个表示，命题4得证。

命题4证明了把 \(E_{f}\) 称为“f 的泛分解代数”这一名称的合理性。关系式 \(f(\mathbf{X}) = \prod_{i=1}^{n} (\mathbf{X} - x_{i,f})\) 称为“f 的泛分解”。设 \(\sigma \in S_{n}\) 为一个置换；由于 \(f(\mathbf{X}) = \prod_{i=1}^{n} (\mathbf{X} - x_{\sigma(i),f})\) ，存在一个自同构 \(t_{\sigma}\) （A-代数 \(E_{f}\) 的），其特征为 \(t_{\sigma}(x_{i,f}) = x_{\sigma(i),f}\) （ \(1 \leqslant i \leqslant n\) ）。我们有 \(t_{\sigma\tau} = t_{\sigma} \circ t_{\tau}\) （对 \(\sigma, \tau\) ，均在 \(S_{n}\) 中），从而得到群 \(S_{n}\) 在 A-代数 \(E_{f}\) 上的作用。

命题5. --- 在泛分解代数 \(E_{f}\) 中，单项式族 \(x_{1}^{\nu(1)}\ldots x_{n}^{\nu(n)}\) （其中 \(0 \leqslant \nu(i) < i\) ， \(1 \leqslant i \leqslant n\) ）是 A-模 \(E_{f}\) 的一组基。特别地， \(E_{f}\) 是秩为 \(n!\)

的自由 A-模。记 \(\mathbf{B} = \mathbf{A}[X_1, \ldots, X_n]\) ，且 \(\mathbf{C} = \mathbf{A}[X_1, \ldots, X_n]^{\mathrm{sym}}\) 。由定理1（IV，第62页），我们有 \(\mathbf{C} = \mathbf{A}[s_1, \ldots, s_n]\) ，且 \(s_1, \ldots, s_n\) 在 \(\mathbf{A}\) 上代数无关。 \(s_1, \ldots, s_n\) 中无常数项的多项式构成一个理想 \(\mathbf{C}^+\) （在 \(\mathbf{C}\) 中），它是 \(\mathbf{A}\) 的补，且由 \(s_1, \ldots, s_n\) 生成。设 \(\mathfrak{c}\) 为 \(\mathbf{C}\) 的由 \(s_1 + a_1\) , \(s_2 - a_2\) , \ldots, \(s_n + (-1)^{n+1}a_n\) 生成的理想。存在一个 \(\mathbf{A}\) -代数自同构（作用于 \(\mathbf{C}\) ），它把 \(s_k\) 映为 \(s_k + (-1)^{k+1}a_k\) （ \(1 \leqslant k \leqslant n\) ），因而把 \(\mathbf{C}^+\) 映到 \(\mathfrak{c}\) 上；因此我们有 \(\mathbf{C} = \mathbf{A} \oplus \mathfrak{c}\) 。此外，IV第62页定理1的 c) 表明

\[
\mathbf {B} = \oplus \mathbf {C X} ^ {\nu}
\]

，其中 S 是由所有满足下述条件的 \(\nu \in \mathbf{N}^n\) 组成的集合： \(0 \leqslant \nu(i) < i\) （ \(1 \leqslant i \leqslant n\) ）。B 的理想 \(\mathfrak{a}\) 由 \(\mathfrak{c}\) 生成，从而 \(\mathfrak{a} = \mathbf{B}\mathfrak{c} = \bigoplus_{\nu \in \mathbb{S}} \mathfrak{c} \cdot \mathbf{X}^\nu\) 。由于 \(\mathbb{C} = \mathbb{A} \oplus \mathfrak{c}\) ，我们得到

\[
\mathbf {B} = \mathfrak {a} \oplus \bigoplus_ {\nu \in S} \mathbf {A X} ^ {\nu},
\]

由此得出命题5，因为 \(\mathbf{E}_f = \mathbf{B} / \mathfrak{a}\) 。

推论。--- A 到首一多项式 \(f \in A[X]\) 的泛分解代数中的典范同态是单射的。

因为 \(E_{f}\) 的单位元构成 A-模 \(E_{f}\) 的一组基的一部分。

命题6. --- 设 \(f \in \mathbf{A}[\mathbf{X}]\) 是一个次数为 \(n\) 的首一多项式，并设 \(P\) 是 \(\mathbf{X}_1, \ldots, \mathbf{X}_n\) 的一个对称多项式，其系数在 \(A\) 中。则恰好存在一个元素 \(a\) （属于 \(A\) ），它具有以下性质：

(FS) 对任意环同态 \(\rho: \mathbf{A} \to \mathbf{B}\) 以及分解 \({}^{\rho} f(\mathbf{X}) = \prod_{i=1}^{n} (\mathbf{X} - \xi_i)\) （在 \(\mathbf{B}[\mathbf{X}]\) 中），我们有 \(\rho(a) = \mathrm{P}(\xi_1, \ldots, \xi_n)\) 。

记 \(f = \mathbf{X}^n + \sum_{k=1}^{n} a_k \mathbf{X}^{n-k}\) ，则由IV第62页的定理1，存在一个多项式

\(\Pi\) （在 \(n\) 个不定元中，系数在 \(\mathbf{A}\) 中），使得 \(\mathbf{P} = \Pi(s_1, \ldots, s_n)\) 。令 \(a = \Pi(-a_1, a_2, \ldots, (-1)^n a_n)\) 。在假设(FS)下，我们有

\[
s _ {k} (\xi_ {1}, \dots , \xi_ {n}) = (- 1) ^ {k} \rho (a _ {k})
\]

由此

\[
\begin{array}{r l} \rho (a) & = \Pi (- \rho (a _ {1}), \rho (a _ {2}), \dots , (- 1) ^ {n} \rho (a _ {n})) \\ & = \Pi (s _ {1} (\xi_ {1}, \dots , \xi_ {n}), \dots , s _ {n} (\xi_ {1}, \dots , \xi_ {n})) \\ & = P (\xi_ {1}, \dots , \xi_ {n}). \end{array}
\]

这证明了满足(FS)的元素 a 的存在性。a 的唯一性由命题5的推论得出，因为我们有关系 \(a \cdot 1 = P(x_{1,f}, \ldots, x_{n,f})\) （在泛分解代数 \(\mathbf{E}_f\) 中）。

采用命题6的记号，有时我们写 \(a = \mathrm{P}^{\#}(f)\) 。下面是一些例子。

例子。--- 1) 如果 \(\mathrm{P} = s_k\) ，则 \(\mathrm{P}^{\#}(f) = (-1)^k a_k\) 。

2) 设 \(g\) 是 \(\mathbf{A}[\mathbf{X}]\) 中的一个多项式，并令

\[
\mathrm{P} \left(\mathrm{X} _ {1}, \dots , \mathrm{X} _ {n}\right) = g \left(\mathrm{X} _ {1}\right) \dots g \left(\mathrm{X} _ {n}\right).
\]

则 \(\mathrm{P}^{\#}(f)\) 正是结式 \(\operatorname{res}(f, g)\) ，由IV第80页推论1。

\begin{enumerate}
\def\labelenumi{\arabic{enumi})}
\setcounter{enumi}{2}
\item
  令 \(\Delta(X_1, \ldots, X_n) = \prod_{i < j} (X_i - X_j)^2\) ，则 \(\Delta^{\#}(f)\) 正是首一多项式 \(f\) 的判别式（IV，第82页，公式（46））。
\item
  令 \(\mathrm{P}(\mathrm{X}_{1},\ldots,\mathrm{X}_{n})=\mathrm{X}_{1}^{k}+\cdots+\mathrm{X}_{n}^{k}\) ；此外，定义代数 \(\mathrm{E}=\mathrm{A}[\mathrm{X}]/(f)\) ，并用 x 表示 X 在 E 中的像。我们回顾一下，A-模 E 是自由的，其基为 \((1,x,\ldots,x^{n-1})\) （IV，第11页，推论）。让我们证明
\end{enumerate}

\[
\operatorname{Tr} _ {\mathrm{E/A}} (x ^ {k}) = \mathrm{P} ^ {\#} (f).\tag{25}
\]

令 \(\pi_k = \mathrm{Tr}_{\mathrm{E / A}}(x^k)\) （对每个整数 \(k \geqslant 1\) ）。考虑到牛顿关系式（IV，第70页），我们只需建立以下关系式

\begin{enumerate}
\def\labelenumi{(\arabic{enumi})}
\setcounter{enumi}{25}
\tightlist
\item
\end{enumerate}

\[
\begin{array}{l l} \pi_ {k} + a _ {1} \pi_ {k - 1} + \dots + a _ {k - 1} \pi_ {1} + k a _ {k} = 0 & \text { for } \quad 1 \leqslant k \leqslant n \\ \pi_ {k} + a _ {1} \pi_ {k - 1} + \dots + a _ {n - 1} \pi_ {k - n + 1} + a _ {n} \pi_ {k - n} = 0 & \text { for } \quad k > n \end{array}\tag{27}
\]

（我们也把它们称为“牛顿关系式”）。关系式(27)是显然的，因为其左边是以下元素的迹

\[
x ^ {k - n} \left(x ^ {n} + a _ {1} x ^ {n - 1} + \dots + a _ {n - 1} x + a _ {n}\right) = 0.
\]

假设 \(1 \leqslant k \leqslant n\) ，并令

\[
y = x ^ {k} + a _ {1} x ^ {k - 1} + \dots + a _ {k - 1} x + a _ {k} \cdot 1;
\]

，令 \(\mathbf{M} = (m_{ij})\) 为线性映射 \(u\mapsto yu\) （在 E 中）关于基 \((x^{i})_{0\leqslant i\leqslant n-1}\) 的矩阵。我们容易得到关系式

\[
\begin{array}{l l} m _ {i i} = a _ {k} & \text { for } \quad 0 \leqslant i <   n - k \\ m _ {i i} = 0 & \text { for } \quad n - k \leqslant i <   n, \end{array}
\]

由此

\[
\operatorname{Tr} _ {\mathrm{E/A}} (y) = \sum_ {i = 0} ^ {n - 1} m _ {i i} = (n - k) a _ {k}.
\]

此外，我们有

\[
\operatorname{Tr} _ {\mathrm{E/A}} (y) = \pi_ {k} + a _ {1} \pi_ {k - 1} + \dots + a _ {k - 1} \pi_ {1} + n a _ {k},
\]

由此得出公式(26)。

